\par In this chapter, we consider modal languages that only use modal terms in which every connective, other than negation,  is expressible in terms of modalities and negation, and all modalities have a uniform relational semantics parametric in a modal term. Moreover, modalities can be composed to form compound modalities, the relational semantics of which are given naturally and uniformly in terms of the composition of the accessibility relations corresponding to the modalities being composed. These languages are called pure modal languages. The language used in this thesis is adapted from~\cite{goranko2006elementary} and is discussed in Section~\ref{pure modal languages section}. All other sections report original research, unless stated otherwise. Section~\ref{pure tableau section} introduces the tableau system $\pmt$ that we will be using, then we do a couple of examples in Section~\ref{pure tableau examples section} to illustrate reasoning in purely modal languages. After this, we prove the soundness and completeness of $\pmt$ in Sections~\ref{soundness section} and Section~\ref{completeness section}, respectively.